\section{分布与离散信号的谱密度}\label{sec:Discrete_Distribution}

如何分析$l^1$以外的离散信号，例如$x[n]=\cos(\Omega_0 n)$的频率成分？

一种直接的想法是利用z变换。然而，\refsub{dtft}中指出，离散信号频谱在z变换的收敛域包含单位元的情况下就是z变换在单位圆上的取值，当单位圆不在z变换的收敛域中时，这一方法就失效了——对于$l^1$以外的信号，z变换仅仅是一种变换域分析的手段，无法展示信号的频率成分。例如，对于$x[n]=\cos(\Omega_0 n)$，根据\cref{def:ROC}，$r=R=1$，$x[n]$的z变换没有收敛域。

可见，目前为止，我们还没有一种合适的工具来分析此类信号，包括最基本的离散复指数信号和离散正余弦信号。直观上，我们需要引入冲击项$\delta$来描述这些信号的频谱，需要做的工作与\refsec{Distributions}是类似的。

\subsection{离散时间傅里叶变换的分布理论}\label{sub:DTFT_distribution}

\begin{definition}[测试函数]
    在DTFT中，测试函数空间$\mathcal{C}$为周期为$2\pi$的光滑函数（即无限阶可导的函数），其中的元素称为\textbf{测试函数}。
\end{definition}

\begin{definition}[速降离散信号]
    如果离散信号$\phi[n]$的衰减速度比任意多项式更快：
    $$\lim_{|n|\to\infty}|n|^k\phi[n]=0,\quad \forall k\in\mathbb{N}$$
    则称$\phi[n]$是速降离散信号，将速降离散信号空间记为$\mathcal{S}$.
\end{definition}

\begin{proposition}
    设$\Phi(\Omega)$是周期为$2\pi$的函数，其傅里叶系数记为$\phi[n]$，则：
    $$\phi[n]\in\mathcal{S}\iff \Phi(\Omega)\in\mathcal{C}$$
\end{proposition}
证明此命题需要用到附录中的结论。
\begin{proof}
    一方面，当$\phi[n]\in\mathcal{S}$时，DTFT满足
    \begin{align*}
        \left|\Phi(\Omega)\right|=\left|\sum_{n=-\infty}^\infty \phi[n]\e^{\im\Omega n}\right|\le \sum_{n=-\infty}^\infty |\phi[n]|
    \end{align*}
    容易由$\phi[n]\in\mathcal{S}$的定义得到$\sum_{n=-\infty}^\infty |\phi[n]|<\infty$，根据魏尔斯特拉斯M判别法，$\sum_{n=-\infty}^\infty \phi[n]\e^{\im\Omega n}$是一致收敛的函数项级数，逐项求导所得的函数项级数为：
    \begin{align*}
        \sum_{n=-\infty}^\infty \phi[n]\left(\e^{\im\Omega n}\right)'=\im\sum_{n=-\infty}^\infty n\phi[n]\e^{\im\Omega n}
    \end{align*}
    类似可得：
    $$\phi[n]\in\mathcal{S}\implies n\phi[n]\in\mathcal{S}\implies \im\sum_{n=-\infty}^\infty n\phi[n]\e^{\im\Omega n}\text{一致收敛}$$
    因此$\Phi'(\Omega)$可导，并且
    $$\Phi'(\Omega)=\im\sum_{n=-\infty}^\infty n\phi[n]\e^{\im\Omega n}$$
    可以归纳地重复上述过程，最终得到$\Phi(\Omega)\in\mathcal{C}$.
    
    另一方面，根据\cref{thm:converge_speed}，$\Phi(\Omega)\in\mathcal{C}\implies\phi[n]\in\mathcal{S}$.
\end{proof}

\begin{definition}[缓增离散信号]
    如果存在常数$C$和$k\in\mathbb{N}$，使得
    $$|x[n]|\le C(1+|n|)^k,\quad\forall n\in\mathbb{Z}$$
    则称$x[n]$是\textbf{缓增离散信号}，将缓增离散信号空间记为$\mathcal{T}$.
\end{definition}

使用该定义描述缓增离散信号是合理的：$(1+|n|)^k$中，1放宽了对信号在$|n|$较小时的要求，而$|n|^k$描述离散信号递增速度不超过$k$次的多项式。

$\mathcal{C}$的对偶空间是分布空间$\mathcal{D}$，速降离散信号空间$\mathcal{S}$的对偶空间是缓增离散信号空间$\mathcal{T}$.

下面定义时域和频域信号的配对。设$x[n],\phi[n]\in l^1\cap l^2$，$X(\e^{\im\Omega})$为$x[n]$的DTFT，$\phi[n]$为$\Phi(\Omega)$的傅里叶系数，则
$$\IDTFT(\Phi(\Omega))[n]=\frac{1}{2\pi}\int_{-\pi}^\pi \Phi(\Omega)\e^{\im\Omega n}\,d\Omega=\phi[-n]$$
在\cref{thm:DTFT_plancherel}中取$y[n]=\overline{\phi[n]}$，则$Y(\e^{\im\Omega})=\overline{\Phi(\Omega)}$，得到：
$$\sum_{n=-\infty}^\infty x[n]\phi[n]=\frac{1}{2\pi}\int_{-\pi}^\pi X(\e^{\im\Omega})\Phi(\Omega)\,d\Omega$$

对于分布做同样的定义：

\begin{definition}[配对]
    时域信号的\textbf{配对}(pairing)定义为：
    $$\langle x,\phi\rangle=\sum_{n=-\infty}^\infty x[n]\phi[n],\quad x\in \mathcal{T},\phi\in\mathcal{S}$$
    其中$x[n]$为缓增离散信号而$\phi[n]$为速降离散信号。

    数字角频域信号的配对定义为：
    $$\langle X,\Phi\rangle=\frac{1}{2\pi}\int_{-\pi}^\pi X(\e^{\im\Omega})\Phi(\Omega)\,d\Omega,\quad X\in\mathcal{D},\Phi\in\mathcal{C}$$
\end{definition}

\begin{definition}[缓增离散信号的离散时间傅里叶变换]
    对于缓增离散信号$x[n]$，称分布$X$是$x[n]$的DTFT，如果：
    $$\langle X,\Phi\rangle=\langle x,\phi\rangle,\quad \forall \phi\in\mathcal{S}$$
\end{definition}

\begin{example}[复指数信号的离散时间傅里叶变换]\label{exmp:exp_DTFT}
    求复指数信号$x[n]=\e^{\im\Omega_0 n}(-\pi<\Omega_0\le\pi)$的DTFT.

    解：记$x[n]$的DTFT为$X$，由定义，
    $$\langle X,\Phi\rangle=\langle x,\phi\rangle=\sum_{n=-\infty}^{\infty}\e^{\im\Omega_0 n}\phi[n],\quad \forall \phi\in\mathcal{S}$$
    光滑函数的傅里叶级数一致收敛于自身（\cref{thm:uniform_convergency}），因此
    $$\Phi(\Omega)=\sum_{n=-\infty}^\infty \phi[n]\e^{\im\Omega n}$$
    可见：（注意数字角频域配对中的系数$2\pi$）
    $$\langle X,\Phi\rangle=\sum_{n=-\infty}^{\infty}\e^{\im\Omega_0 n}\phi[n]=\Phi(\e^{\im\Omega_0})=\langle 2\pi\delta_{\Omega_0},\Phi\rangle$$
    数字角频域总是以$2\pi$为周期的，根据以往的经验，$x[n]$的DTFT应该是以下周期分布：
    $$X=2\pi\sum_{k=-\infty}^{\infty}\delta_{\Omega_0+2\pi k}=2\pi\tau_{\Omega_0}\shah_{2\pi}$$
    其中，平移算子$\tau$的定义可以参考\cref{def:distribution_delay}，$\shah_{2\pi}$是\cref{def:shah}给出的以$2\pi$为周期的狄拉克梳状分布。

    但是，事实上，我们不需要采用像$2\pi\tau_{\Omega_0}\shah_{2\pi}$这么复杂的符号，这一方面是因为$2\pi\delta_{\Omega_0}$已经足以描述$X$的行为，并且其意义是不言自明的；另一方面，从本节中对于配对的定义可以看出，$\mathcal{C,D}$是否进行周期化并不影响分布的运算，换言之，$2\pi\delta_{\Omega_0}$和$2\pi\tau_{\Omega_0}\shah_{2\pi}$在$\mathcal{D}$中\textbf{完全是同一个分布}。因此，$X=2\pi\delta_{\Omega_0}$.
\end{example}

作为$\Omega_0=0$时的特例，常数序列$x[n]=1$的DTFT为$2\pi\delta$.

在开始讨论分布理论中的DTFT的性质之前，我们仿照\refsec{Distributions}给出以下定义。

\begin{definition}[分布与函数的乘法]
    设$X\in\mathcal{D},g,\Phi\in\mathcal{C}$，则$g$与$X$的乘积$gX$由下式定义：
    $$\langle gX,\Phi\rangle=\langle X,g\Phi\rangle$$
\end{definition}

\begin{definition}[分布的微分]
    设$X\in\mathcal{D},\Phi\in\mathcal{C}$，则$X$的微分$X'$由下式定义：
    $$\langle X',\Phi\rangle=-\langle X,\Phi'\rangle$$
\end{definition}

\begin{definition}[分布的平移与伸缩]
    设$X\in\mathcal{D},\Phi\in\mathcal{C}$，则$X$向右平移$\Omega_0$所得分布记为$\tau_{\Omega_0}X$，满足：
    $$\langle \tau_{\Omega_0}X,\Phi \rangle=\langle X,\tau_{-\Omega_0}\Phi \rangle$$
    $\sigma_a X$由下式定义：
    $$\langle \sigma_a X,\Phi \rangle=\left\langle X,\frac{1}{|a|}\sigma_{1/a}\Phi \right\rangle$$
    特别地，$a=-1$时，$\sigma_{a}X$就是$X$的翻转，记作$X^-$.$\langle X^-,\Phi \rangle=\left\langle X,\Phi^- \right\rangle$.
\end{definition}

\begin{definition}[分布的共轭]
    设$X\in\mathcal{D},\Phi\in\mathcal{C}$，则$X$的共轭$\overline{X}$由下式定义：
    $$\langle \overline{X},\Phi\rangle=\overline{\langle X,\overline{\Phi}\rangle}$$
    由此可以定义分布的实部与虚部：
    $$\mathrm{Re}(X)=\frac{X+\overline{X}}{2},\quad \mathrm{Im}(X)=\frac{X-\overline{X}}{2\im}$$
\end{definition}

\begin{definition}[分布的卷积]
    设$X,Y\in\mathcal{D},\Phi\in\mathcal{C}$，则$X,Y$的周期卷积为
    $$\langle X*Y,\Phi\rangle=\langle X(u),\langle Y(v),\Phi(u+v)\rangle\rangle$$
\end{definition}
与\cref{def:conv_distr_distr}不同，由于是有限区间上的积分，分布的卷积总是有定义的。%与\cref{exmp:associative_loss}的比较？

请读者自行验证上述定义对于正则分布成立。周期函数的卷积见\ref{def:period_conv}。

容易验证，对于时域的配对而言，同样有：
$$\langle gx,\phi\rangle=\langle x,g\phi\rangle,\quad \langle \tau_{n_0}x,\phi \rangle=\langle x,\tau_{-n_0}\phi \rangle,\quad \langle \overline{x},\phi\rangle=\overline{\langle x,\overline{\phi}\rangle},\quad x,\phi,g\in\mathcal{S},n_0\in\mathbb{Z}$$

此外，设$x\in\mathcal{T},\phi\in\mathcal{S}$，考虑将时域卷积写成配对的形式：
$$\langle x*y,\phi\rangle=\langle x,(y^-)*\phi\rangle$$
此时我们希望$$x*y\in\mathcal{T},\quad (y^-)*\phi\in\mathcal{S}$$
与模拟信号的情况一样，可以有多种条件用来保证这一点，其中一种条件为$y\in\mathcal{S}$，这也是我们之后主要考虑的情况。

\subsection{离散时间傅里叶变换的分布的运算性质}\label{sub:DTFT_distribution_operation}

本小节中依次验证\refsub{property_of_dtft}中的DTFT的性质。不做特殊说明时默认:
$$x[n],y[n]\in\mathcal{T},\quad \phi[n]\in\mathcal{S}$$
并用符号$X,Y\in\mathcal{D}$表示$x,y$的DTFT，$\phi[n]$是$\Phi(\Omega)$的傅里叶系数，即用$\Phi(\Omega)\in\mathcal{C}$表示$\phi^-$的DTFT.

\begin{proposition}[线性性]
    $$\DTFT(ax[n]+by[n])=aX+bY,\quad\forall a,b\in\mathbb{C}$$
\end{proposition}
该命题直接得自配对$\langle\cdot,\cdot\rangle$对两个变元的线性性。

\begin{proposition}[时移性质]
    $$\DTFT(x[n-n_0])=\e^{-\im \Omega n_0}X$$
\end{proposition}
\begin{proof}
    $$\phi[n+n_0]=\frac{1}{2\pi}\int_{-\pi}^\pi \Phi(\Omega)\e^{-\im\Omega (n+n_0)}\,d\Omega$$
    因此$\phi[n+n_0]$是$\Phi(\Omega)\e^{-\im\Omega n_0}$的傅里叶系数。
    \begin{align*}
        \langle \tau_{n_0}x,\phi\rangle&=\langle x,\tau_{-n_0}\phi\rangle=\langle X,\e^{-\im\Omega n_0}\Phi\rangle=\langle\e^{-\im \Omega n_0}X,\Phi\rangle
    \end{align*}
    即$\DTFT(x[n-n_0])=\e^{-\im \Omega n_0}X$.
\end{proof}

\begin{proposition}[频移性质]
    $$\DTFT(\e^{\im\Omega_0 n}x[n])=\tau_{\Omega_0}X$$
\end{proposition}
\begin{proof}
    $$\e^{\im\Omega_0 n}\phi[n]=\frac{1}{2\pi}\int_{-\pi}^\pi \Phi(\Omega)\e^{-\im(\Omega-\Omega_0) n}\,d\Omega=\frac{1}{2\pi}\int_{-\pi}^\pi \Phi(\Omega+\Omega_0)\e^{-\im\Omega n}\,d\Omega$$
    因此$\e^{\im\Omega_0 n}\phi[n]$是$\Phi(\Omega+\Omega_0)$的傅里叶系数。
    \begin{align*}
        \langle \e^{\im\Omega_0 n}x,\phi\rangle&=\langle x,\e^{\im\Omega_0 n}\phi\rangle=\langle X,\tau_{-\Omega_0}\Phi\rangle=\langle \tau_{\Omega_0}X,\Phi\rangle
    \end{align*}
    即$\DTFT(\e^{\im\Omega_0 n}x[n])=\tau_{\Omega_0}X$.
\end{proof}

\begin{proposition}[共轭性质]
    $$\DTFT\left(\overline{x[n]}\right)=\overline{X^-}$$
\end{proposition}
\begin{proof}
    $$\overline{\phi[n]}=\frac{1}{2\pi}\overline{\int_{-\pi}^\pi \Phi(\Omega)\e^{-\im\Omega n}\,d\Omega}=\frac{1}{2\pi}\int_{-\pi}^\pi \overline{\Phi(\Omega)}\e^{\im\Omega n}\,d\Omega=\frac{1}{2\pi}\int_{-\pi}^\pi \overline{\Phi(-\Omega)}\e^{-\im\Omega n}\,d\Omega$$
    因此$\overline{\phi[n]}$是$\overline{\Phi(-\Omega)}$的傅里叶系数。
    \begin{align*}
        \langle \overline{x},\phi\rangle&=\overline{\langle x,\overline{\phi}\rangle}=\overline{\langle X,\overline{\Phi^-}\rangle}=\langle \overline{X^-},\Phi\rangle
    \end{align*}
    最后一个等号用到了$\overline{X}^-=\overline{X^-}$.
\end{proof}
\begin{corollary}
    如果$x[n]$为实数序列，则$X$具有共轭对称性：$X=\overline{X^-}$.
\end{corollary}

\begin{proposition}[反转性质]
    $$\DTFT(x^-[n])=X^-$$
\end{proposition}
\begin{proof}
    $$\phi[-n]=\frac{1}{2\pi}\int_{-\pi}^\pi \Phi(\Omega)\e^{\im\Omega n}\,d\Omega=\frac{1}{2\pi}\int_{-\pi}^\pi \Phi(-\Omega)\e^{-\im\Omega n}\,d\Omega$$
    因此$\phi[-n]$是$\Phi(-\Omega)$的傅里叶系数。
    \begin{align*}
        \langle x^-,\phi\rangle=\langle x,\phi^-\rangle=\langle X,\Phi^-\rangle=\langle X^-,\Phi\rangle
    \end{align*}
\end{proof}

\begin{proposition}[频域微分性质]\label{prop:DTFT_derivative}
    $$\DTFT(nx[n])=\im X'$$
\end{proposition}
\begin{proof}
    $\im \Phi'$的傅里叶系数为
    \begin{align*}
        \frac{1}{2\pi}\int_{-\pi}^\pi \im \Phi'(\Omega)\e^{-\im \Omega n}\,d\Omega&=\frac{\im}{2\pi}\left(\evalat{\Phi(\Omega)\e^{-\im \Omega n}}{-\pi}{\pi}-\int_{-\pi}^\pi \Phi(\Omega)(\e^{-\im \Omega n})'\,d\Omega\right)\\
        &=\frac{n}{2\pi}\int_{-\pi}^\pi \Phi(\Omega)\e^{-\im\Omega n}\,d\Omega\\
        &=n\phi[n]
    \end{align*}
    因此
    \begin{align*}
        \langle nx,\phi\rangle=\langle x,n\phi\rangle=\langle X,\im\Phi'\rangle=\langle \im X',\Phi\rangle
    \end{align*}
\end{proof}

\begin{theorem}[时域卷积定理]
    设$x\in\mathcal{T},y,\phi\in\mathcal{S}$，则
    $$\DTFT(x*y)=XY$$
\end{theorem}
\begin{proof}
    \begin{align*}
        \langle(x*y),\phi\rangle=\langle x,(y^-)*\phi\rangle
    \end{align*}
    令$\psi[n]=(y^-)*\phi[n]$，对应的频域函数为
    \begin{align*}
        \Psi(\Omega)&=\sum_{k=-\infty}^\infty \psi[n]\e^{\im\Omega k}\\
        &=\sum_{k=-\infty}^\infty \left(\sum_{m=-\infty}^\infty y[m]\phi[k+m]\right)\e^{\im\Omega k}\\
        &=\sum_{m=-\infty}^\infty y[m]\e^{-\im\Omega m}\sum_{n=-\infty}^\infty \phi[n]\e^{\im \Omega n}&(n=k+m)\\
        &=Y(\e^{\im\Omega})\Phi(\Omega)
    \end{align*}
    因此\begin{align*}
        \langle(x*y),\phi\rangle=\langle x,(y^-)*\phi\rangle=\langle X,Y\Phi\rangle=\langle XY,\Phi\rangle
    \end{align*}
\end{proof}

\begin{theorem}[频域卷积定理]
    $$\DTFT(x[n]y[n])=X*Y$$
\end{theorem}
\begin{proof}
    \begin{align*}
        \langle X*Y,\Phi\rangle&=\langle X(u),\langle Y(v),\Phi(u+v)\rangle\rangle
    \end{align*}
    令$\Psi(u)=\langle Y(v),\Phi(u+v)\rangle$，则其傅里叶系数为
    \begin{align*}
        \psi[n]&=\frac{1}{2\pi}\int_{-\pi}^\pi \Psi(u)\e^{-\im un}\,du\\
        &=\left\langle Y(v),\frac{1}{2\pi}\int_{-\pi}^\pi \Phi(u+v)\e^{-\im un}\,du\right\rangle\\
        &=\left\langle Y(v),\frac{1}{2\pi}\int_{-\pi}^\pi \Phi(\Omega)\e^{-\im (\Omega-v)n}\,d\Omega\right\rangle&(\Omega=u+v)\\
        &=\left\langle Y(v),\phi[n]\e^{\im vn}\right\rangle\\
        &=\frac{1}{2\pi}\phi[n]\int_{-\pi}^\pi Y(v)\e^{\im vn}\,dv\\
        &=\phi[n]y[n]
    \end{align*}
    因此
    \begin{align*}
        \langle X*Y,\Phi\rangle&=\langle X(u),\langle Y(v),\Phi(u+v)\rangle\rangle\\
        &=\langle X,\Psi\rangle\\
        &=\langle x,\phi y\rangle\\
        &=\langle xy,\phi\rangle
    \end{align*}
\end{proof}

由于$x[n],y[n]\in\mathcal{T}$时，$x[n]y[n]\in\mathcal{T}$，$X*Y$总是存在的，且\textbf{满足交换律和结合律}。

\begin{example}[离散正弦、余弦函数的离散时间傅里叶变换]\label{exmp:cos_DTFT}
    根据欧拉公式及DTFT的线性性可得：
    $$\DTFT(\cos(\Omega_0 n))=\pi(\delta_{\Omega_0}+\delta_{-\Omega_0}),\quad \DTFT(\sin(\Omega_0 n))=-\pi\im(\delta_{\Omega_0}-\delta_{-\Omega_0})$$
    根据\cref{prop:discrete_period}，当$\Omega_0/\pi$是无理数时，$\cos(\Omega_0 n)$和$\sin(\Omega_0 n)$不具有周期性，但是总存在有理数与$\Omega_0/\pi$充分靠近，并且频谱为线谱，因此可以说这种信号是\textbf{几乎周期}(almost periodic)的。我们在\refsub{power_signal}中讨论过类似的概念。
\end{example}

在上例中，令$\Omega_0=\pi$，则
$$\DTFT(\cos(\pi n))=\DTFT((-1)^n)=\pi(\delta_\pi+\delta_{-\pi})=2\pi\delta_\pi$$

\begin{example}[单位阶跃序列的离散时间傅里叶变换]
    \begin{align*}
        &u[n]-u[n-1]=\delta[n]\\
        \implies&\langle (1-\tau_1)u,\phi\rangle=\langle\delta,\phi\rangle\\
        \implies&\langle(1-\e^{\im\Omega})U,\Phi\rangle=\langle 1,\Phi\rangle
    \end{align*}
    然而，作为一个分布方程，$(1-\e^{\im\Omega})U=1$并不能表明$U=1/(1-\e^{\im\Omega})$，由于
    $$\frac{1}{1-\e^{\im\Omega}}\sim\frac{\im}{\Omega},\quad \Omega\to 0$$
    $1/(1-\e^{\im\Omega})$在$[-\pi,\pi]$上不可积，也就无法与$1\in\mathcal{C}$配对。我们需要更加仔细地考察$U$的性质。

    $(1-\e^{\im\Omega})U=1$确定了$U$在$\Omega\neq 0$处的值，为了使它能与1配对，我们只能认为$\langle U,\Phi\rangle$的意义是在主值意义下的积分，这样，$U$仍满足它的分布方程：
    $$(1-\e^{\im\Omega})U=(1-\e^{\im\Omega})\pv\frac{1}{(1-\e^{\im\Omega})}=1$$
    由于$\Phi\in\mathcal{C}$光滑，$1/(1-\e^{\im\Omega})$在$\Omega=0$附近的表现类似于奇函数$\im/\Omega$，该主值意义下的积分总是收敛的。

    主值积分不能描述$\Omega=0$处的$U$的行为，由于它是$1/(1-\e^{\im\Omega})$的极点，直观上可以使用$\delta$分布来处理这个问题。我们引入另一个分布方程：
    $$u[n]+u[-n]=\begin{cases}
        2,&\text{if }n=0\\
        1,&\text{if }n\neq 0
    \end{cases}=1+\delta[n]\implies U+U^-=2\pi\delta+1$$
    代入
    $$U=\pv\dfrac{1}{1-\e^{\im\Omega}}+a\delta$$
    得到$a=\pi$，因此：
    $$U=\pv\dfrac{1}{1-\e^{\im\Omega}}+\pi\delta$$
    对于z变换的收敛域边界恰为单位圆的信号，这样的作法是普适的。
\end{example}
\begin{example}
    本例探究$\delta'$的性质。
    $$\langle \delta',\Phi\rangle=-\langle\delta,\Phi'\rangle=-\Phi'(0)$$
    这与\cref{exmp:delta'}中的结论是类似的。根据\cref{prop:DTFT_derivative}，有：
    $$\DTFT(n)=\im\delta'$$
    这表明$\delta'$是$x[n]=-\im n$的DTFT.
\end{example}

利用分布的卷积的定义，容易得到
$$\delta_a*\delta_b=\delta_{a+b},\quad \delta_a*X=\tau_a X$$
基于此，利用频域卷积定理可得：（注意周期函数的卷积带有系数$1/2\pi$）

\begin{example}[单边正弦、余弦序列的离散时间傅里叶变换]
    \begin{align*}
        \DTFT(\cos(\Omega_0 n)u[n])&=\pi(\delta_{\Omega_0}+\delta_{-\Omega_0})*U\\
        &=\pv\frac{1}{2}\left(\frac{1}{1-\e^{\im(\Omega-\Omega_0)}}+\frac{1}{1-\e^{\im(\Omega+\Omega_0)}}\right)+\frac{\pi}{2}(\delta_{\Omega_0}+\delta_{-\Omega_0})\\
        &=\pv\frac{1-\e^{\im\Omega}\cos\Omega_0}{1-2\e^{\im\Omega}\cos\Omega_0+\e^{2\im\Omega}}+\frac{\pi}{2}(\delta_{\Omega_0}+\delta_{-\Omega_0})\\
        \DTFT(\sin(\Omega_0 n)u[n])&=-\pi\im(\delta_{\Omega_0}-\delta_{-\Omega_0})*U\\
        &=-\frac{\im}{2}\pv\left(\frac{1}{1-\e^{\im(\Omega-\Omega_0)}}-\frac{1}{1-\e^{\im(\Omega+\Omega_0)}}\right)-\frac{\pi\im}{2}(\delta_{\Omega_0}-\delta_{-\Omega_0})\\
        &=-\pv\frac{\e^{\im\Omega}\sin\Omega_0}{1-2\e^{\im\Omega}\cos\Omega_0+\e^{2\im\Omega}}-\frac{\pi\im}{2}(\delta_{\Omega_0}-\delta_{-\Omega_0})
    \end{align*}
    与\cref{exmp:causual_cos_z}和\cref{exmp:casual_sin_z}对比，可以发现这两个序列的DTFT仍是z变换在单位圆上的取值，尽管单位圆并不它们的z变换的收敛域内。
\end{example}

\subsection{分布的离散时间傅里叶变换的连续性}\label{sub:DTFT_continuous}

类似于\cref{thm:distribution_fourier_continuous}，我们有：
\begin{theorem}\label{thm:DTFT_continuous_1}
    设$\{X_k\}_{k=1}^\infty\subset\mathcal{D},X\in\mathcal{D}$，如果$X_k\to X$，则$x_k[n]\to x[n],\forall n\in\mathbb{Z}$.
\end{theorem}
\begin{proof}
    $$\lim_{k\to\infty}\langle x_k,\phi\rangle=\lim_{k\to\infty}\langle X_k,\Phi\rangle=\langle X,\Phi\rangle=\langle x,\phi\rangle$$
    于是
    $$\delta[n-m]\in\mathcal{S},\forall m\in\mathbb{Z}\implies x_k[n]=x[n],\forall n\in\mathbb{Z}$$
\end{proof}

然而，其逆命题需要额外的条件：
\begin{theorem}\label{thm:DTFT_continuous_2}
    设$\{x_k[n]\}_{k=1}^\infty\subset\mathcal{T},x[n]\in\mathcal{T}$，如果$x_k[n]\to x[n],\forall n\in\mathbb{Z}$，并且$\{x_k[n]\}$具有\textbf{一致缓增性}，即：
    $$\exists C>0,M\in\mathbb{N}^*,\quad |x_k[n]|\le C(1+|n|)^M,\quad \forall k\in\mathbb{N}^*,n\in\mathbb{Z}$$
    则$X_k\to X$;
\end{theorem}
\begin{proof}
    设$\phi[n]\in\mathcal{S}$，$\Phi(\Omega)\in\mathcal{C}$以$\phi[n]$为傅里叶系数，我们说明下式中的极限可以与求和交换：\begin{align*}
        \lim_{k\to\infty}\langle x_k,\phi\rangle=\lim_{k\to\infty}\sum_{n=-\infty}^\infty x_k[n]\phi[n]
    \end{align*}
    构造函数列
    $$\{f_k:\mathbb{R}\to\mathbb{C}\}_{k=1}^\infty,\quad f_k(t)=\sum_{n=-\infty}^\infty x_k[n]\phi[n]\chi_{[n,n+1)}(t)$$
    显然$\{f_k\}$具有逐点极限
    $$f:\mathbb{R}\to\mathbb{C},\quad f(t)=\sum_{n=-\infty}^\infty x[n]\phi[n]\chi_{[n,n+1)}(t)$$
    根据速降离散信号的定义，
    $$\lim_{n\to\infty}Cn^2(1+|n|)^M\phi[n]=0$$
    因此$\{Cn^2(1+|n|)^M\phi[n]\}_{k=1}^\infty$有上界$M'$，于是
    $$\left|Cn^2(1+|n|)^M\phi[n]\right|\le \dfrac{M}{n^2},\quad \sum_{n=-\infty}^\infty \left|Cn^2(1+|n|)^M\phi[n]\right|\le\sum_{n=-\infty}^\infty\frac{M}{n^2}<\infty$$
    由此构造出可积的控制函数：
    $$g(t)=\sum_{n=-\infty}^\infty \left|Cn^2(1+|n|)^M\phi[n]\right|\chi_{[n,n+1)}(t),\quad |f_k(t)|\le g(t),\quad \forall k,\forall t$$
    由控制收敛定理（\cref{thm:DCT}），可得：
    $$\int_{\mathbb{R}}f\,dm=\lim_{k\to\infty}\int_{\mathbb{R}}f_k\,dm$$
    即
    $$\langle x,\phi\rangle=\sum_{n=-\infty}^\infty x[n]\phi[n]=\lim_{k\to\infty}\sum_{n=-\infty}^\infty x_k[n]\phi[n]=\lim_{k\to\infty}\langle x_k,\phi\rangle$$
    因此
    $$\langle X,\Phi\rangle=\langle x,\phi\rangle=\langle x_k,\phi\rangle=\lim_{k\to\infty}\langle X_k,\Phi\rangle$$
\end{proof}

一致缓增性的条件不可以去除：
\begin{example}
    构造缓增序列
    $$x_k[n]=\begin{cases}
        k!,&n=k\\
        0,&n\neq k
    \end{cases}$$
    固定$n$，当$k\to\infty$时，$x_k[n]\to 0$，因此$\{x_k[n]\}$逐点收敛到0序列$x[n]\equiv 0$.

    但是，取速降序列$\phi[n]=\dfrac{1}{n!}$，则
    $$\langle x_k,\phi\rangle=\sum_{n=-\infty}^\infty x_k[n]\phi[n]=1,\quad \forall k\in \mathbb{N}^*$$
    而$\langle X,\Phi\rangle=0$，这表明
    $$\lim_{k\to\infty}\langle X_k,\Phi\rangle\neq \langle X,\Phi\rangle$$
\end{example}

\subsection{离散信号的相关和谱密度}\label{sub:discrete_corr}

仿照连续信号，容易提出离散信号的相关的定义：
\begin{definition}[能量离散信号的相关]
    \label{def:discrete_correlation_energy}
    设$x[n],y[n]\in l^2$，它们的\textbf{互相关}(cross correlation)定义为：
    $$r_{xy}[l]=\sum_{n=-\infty}^\infty x[n+l]\overline{y[n]}$$
    当$x=y$时，称之为$x[n]$的\textbf{自相关}(auto-correlation)并记为$r_x[l]$.
\end{definition}

利用与\cref{prop:cross_correlation}和\cref{thm:winner-khinchin}同样的方法，可以证明：

\begin{proposition}[能量信号的相关的性质]
    设$x,y\in l^2$，则\begin{enumerate}
        \item $r_{xy}=x*\overline{y}^-=\overline{r_{gf}^-}$，特别地，如果$x,y$都是实信号，则
        $$r_{xy}=r^-_{yx}=x*y^-=((x^-)*y)^-$$
        \item $r_{\tau_n x,y}=\tau_n r_xy=r_{x,\tau_{-n}y}$
        \item $r_{xy}\le\|x\|_2\|y\|_2$，特别地，$r_x[l]\le r_x(0)=\|x\|_2^2$
    \end{enumerate}
\end{proposition}

\begin{theorem}[能量信号的维纳-辛钦定理]
    设$x,y\in l^1\cap l^2$，则$\DTFT(R_{xy})=X(\e^{\im\Omega})\overline{Y(\e^{\im\Omega})}$
\end{theorem}

\begin{example}[指数序列的自相关]
    设$x[n]=a^n u[n],0<a<1$.由于$x[n]$是实信号，$r_x[l]$是偶对称的，只需计算$l\ge 0$时的相关：
    \begin{align*}
        r_x[l]&=\sum_{n=-\infty}^\infty x[n+l]x[n]\\
        &=\sum_{n=0}^\infty a^{n+l}a^n\\
        &=a^l\sum_{n=0}^\infty a^{2n}\\
        &=\frac{a^l}{1-a^2}
    \end{align*}
    利用$r_x[-l]=r_x[l]$，得到$r_x[l]=\dfrac{a^{|l|}}{1-a^2}$.
\end{example}

由DTFT的帕塞瓦尔定理（\cref{thm:DTFT_plancherel}），有：
$$E_x=\sum_{n=-\infty}^{\infty}|x[n]|^2=\frac{1}{2\pi}\int_{-\pi}^{\pi}|X(\mathrm{e}^{\mathrm{i}\Omega})|^2\,d\Omega$$

这表明$x\in l^1\cap l^2$的\textbf{能量是分布在各个数字角频率上的}，由此定义：

\begin{definition}[能量谱密度]
    $x\in l^1\cap l^2$的\textbf{能量谱密度}(energy spectral density,ESD)定义为
    $$S_x(\Omega)=\dfrac{1}{2\pi}|\DTFT(x)(\Omega)|^2$$
\end{definition}

\begin{example}[指数序列的能量谱密度]
    设$x[n]=a^n u[n],0<a<1$，则$r_x[l]=\dfrac{a^{|l|}}{1-a^2}$，对其做DTFT，得到：
    \begin{align*}
        |\DTFT(x)(\Omega)|^2&=\DTFT(r_x)(\Omega)\\
        &=\sum_{l=-\infty}^\infty \dfrac{a^{|l|}}{1-a^2}\e^{-\im\Omega l}\\
        &=\frac{1}{1-a^2}\left(\sum_{l=0}^\infty a^l(\e^{-\im\Omega l}+\e^{\im\Omega l})-1\right)\\
        &=\frac{1}{1-a^2}\left(\frac{1}{1-a\e^{-\im\Omega}}+\frac{1}{1-a\e^{\im\Omega}}-1\right)\\
        &=\frac{1}{1-a^2}\left(\frac{2(1-a\cos\Omega)}{1-2a\cos\Omega+a^2}-1\right)\\
        &=\frac{1}{1-2a\cos\Omega+a^2}
    \end{align*}
    因此它的能量谱密度为
    $$S_x(\Omega)=\frac{1}{2\pi}|\DTFT(x)(\Omega)|^2=\frac{1}{2\pi(1-2a\cos\Omega+a^2)}$$
    它的能量谱密度如下图所示：
    \begin{figure}[H]
        \centering
        \includegraphics[width=0.6\textwidth]{exp_ESD}
        \caption{指数序列的能量谱密度}
    \end{figure}
\end{example}

仿照\refsub{power_signal}，可以定义\textbf{离散功率信号空间}并赋予内积：

\begin{definition}[离散功率信号空间]
    能量无限而功率非零有限的信号称为\textbf{离散功率信号}：
    $$\sum_{n=-\infty}^\infty |x[n]|^2=\infty,\quad 0<\lim_{M\to\infty}\frac{1}{2M+1}\sum_{n=-M}^M |x[n]|^2<\infty$$
    我们约定：能量为0的信号也属于能量信号，从而能量信号空间就是$l^2$；功率为0的信号也是功率信号，从而功率信号空间构成线性空间，记为$b^2$.
\end{definition}

离散功率信号空间中的典型信号是周期的或者几乎周期的。

\begin{definition}[离散功率信号空间的内积]
    设$x[n],y[n]\in b^2$，它们的内积定义为
    $$(x,y)_{\pwr}=\lim_{M\to\infty}\frac{1}{2M+1}\sum_{n=-\infty}^\infty x[n]\overline{y[n]}$$
\end{definition}

由这个内积导出的范数$\|\cdot\|_{\pwr}$并不具有同一性，但可以通过引入等价关系来解决这一问题：
\begin{align*}
    x\sim y\iff \|x-y\|_{\pwr}=0
\end{align*}

\begin{definition}[功率离散信号的相关]
    \label{def:discrete_correlation_power}
    设$x[n],y[n]\in b^2$，它们的\textbf{互相关}(cross correlation)定义为：
    $$r_{xy}[l]=\lim_{M\to\infty}\frac{1}{2M+1}\sum_{n=-M}^M x[n+l]\overline{y[n]}$$
    当$x=y$时，称之为$x[n]$的\textbf{自相关}(auto-correlation)并记为$r_x[l]$.
\end{definition}

通过截短信号的的极限来定义功率信号的功率谱密度。我们用
$$x_M[n]=\begin{cases}
    x[n],&\text{if }-M\le n\le M\\
    0,&\text{otherwise}
\end{cases}$$
来表示长为$2M+1$的离散信号的截短，并用$X_M(\Omega)$表示其DTFT。对$x_M[n]$用帕塞瓦尔定理，有：
$$E_{x_M}=\sum_{n=-\infty}^\infty |x_M[n]|^2=\frac{1}{2\pi}\int_{-\pi}^\pi |X_M(\e^{\im\Omega})|^2\,d\Omega$$

这样，在$M\to\infty$时，$E_{x_M}$的时间平均值就是$x[n]$的功率：
$$P_x=\lim_{M\to\infty}\frac{1}{2M+1}\sum_{n=-M}^M |x[n]|^2=\lim_{M\to\infty}\frac{E_{x_M}}{2M+1}=\lim_{M\to\infty}\frac{1}{2\pi}\int_{-\pi}^\pi |X_M(\e^{\im\Omega})|^2\,d\Omega$$

由于离散功率信号$x[n]$的DTFT是作为分布来定义的，在分布的意义下可以交换积分与求和的次序，因此可以将功率写为密度函数的积分的形式：
\begin{definition}[离散功率信号的功率谱密度]
    功率信号$x[n]$的\textbf{功率谱密度}(power spectral density)为
    $$S_x(\Omega)=\frac{1}{2\pi}\lim_{M\to\infty}\frac{|X_M(\e^{\im\Omega})|^2}{2M+1}$$
\end{definition}

仿照\cref{def:cross_energy_power}和\cref{def:cross_energy_power_density}，还可以定义离散信号的\textbf{互能量}以及\textbf{互能量谱密度}，功率信号同理：
\begin{definition}[互能量，互功率与谱密度]
    如果$x[n],y[n]\in l^2$，则它们的\textbf{互能量}定义为
    $$E_{xy}=(x,y)=\sum_{n=-\infty}^\infty x[n]\overline{y[n]}$$
    其中$(\cdot,\cdot)$为希尔伯特空间$l^2$上的内积。如果$x[n],y[n]\in l^1\cap l^2$，则它们的\textbf{互能量谱密度}(cross energy spectral density,CESD)定义为
    $$S_{xy}(\Omega)=\frac{1}{2\pi}X(\e^{\im\Omega})\overline{Y(\e^{\im\Omega})}$$

    如果$x[n],y[n]\in b^2$，则它们的\textbf{互功率}定义为
    $$P_{xy}=(x,y)_{\pwr}=\lim_{M\to\infty}\frac{1}{2M+1}\sum_{n=-M}^M x[n]\overline{y[n]}$$
    它们的\textbf{互功率谱密度}(cross power spectral density,CPSD)定义为
    $$S_{xy}(\Omega)=\frac{1}{2\pi}\lim_{M\to\infty}\frac{1}{2M+1}X_M(\e^{\im\Omega})\overline{Y_M(\e^{\im\Omega})}$$
\end{definition}

仿照\cref{thm:power_winner_khinchin}，我们有：
\begin{theorem}[离散功率信号的维纳-辛钦定理]
    \label{thm:deiscrete_winner_khinchin}
    设$x[n],y[n]\in b^2$，则
    $$\DTFT(r_{xy})(\Omega)=\lim_{M\to\infty}\frac{1}{2M+1}X_M(\e^{\im\Omega})\overline{Y_M(\e^{\im\Omega})}$$
\end{theorem}
\begin{proof}
    定义截短信号所得的能量信号$x_M[n]$和$y_M[n]$的互相关函数为
    $$r_M[l]=r_{x_M,y_M}[l]=\sum_{n=-M}^M x[n]\overline{y[n]}$$
    由维纳-辛钦定理，有：
    $$\DTFT(r_M)(\e^{\im\Omega})=X_M(\e^{\im\Omega})\overline{Y(\e^{\im\Omega})}$$
    等式两边同时除以$2M+1$并与测试函数$\Phi(\Omega)\in\mathcal{C}$配对：
    $$\left\langle \frac{1}{2M+1}\DTFT\left(r_M\right),\Phi\right\rangle=\left\langle \frac{1}{2M+1}X_M(\e^{\im\Omega})\overline{Y(\e^{\im\Omega})},\Phi\right\rangle$$
    令$M\to\infty$，由分布的傅里叶变换的连续性（\cref{thm:distribution_fourier_continuous}），左式等于：
    $$\lim_{M\to\infty}\left\langle \frac{1}{2M+1}\DTFT\left(r_M\right),\Phi\right\rangle=\left\langle\DTFT\left(\lim_{M\to\infty}\frac{r_M}{2M+1}\right),\Phi\right\rangle=\langle\DTFT(r_{xy}),\Phi\rangle$$
    右式等于：
    $$\lim_{M\to\infty}\left\langle \frac{1}{2M+1}X_M(\e^{\im\Omega})\overline{Y(\e^{\im\Omega})},\Phi\right\rangle=\langle S_{xy},\Phi\rangle$$
\end{proof}

这样，功率谱可以通过相关函数的DTFT来计算。

\begin{example}
设离散功率信号 \(x[n]=\cos(\Omega_0 n)\)，其中 \(0<\Omega_0<\pi\) 且 \(\Omega_0\neq 0,\pi\)，截短信号为$x_M[n]$，则其自相关函数为
\[
r_M[l]=\sum_{n=-\infty}^{\infty} x_M[n+l]x_M[n].
\]
由于两个截短信号的支集分别为 \([-M,M]\) 和 \([-M-l,M-l]\)，其交集非空当且仅当 \(|l|\le 2M\)。于是
\[
r_M[l]=
\begin{cases}
\displaystyle \sum_{n=-M}^{M-l}\cos(\Omega_0(n+l))\cos(\Omega_0 n), & 0\le l\le 2M,\\[1.2ex]
\displaystyle \sum_{n=-M-l}^{M}\cos(\Omega_0(n+l))\cos(\Omega_0 n), & -2M\le l<0,\\[1.2ex]
0, & |l|>2M.
\end{cases}
\]
利用积化和差
\[
\cos(\Omega_0(n+l))\cos(\Omega_0 n)
=\frac12\left[\cos(\Omega_0(2n+l))+\cos(\Omega_0 l)\right].
\]
当 \(0\le l\le 2M\) 时，
\[
r_M[l]=\frac12\left[
\sum_{n=-M}^{M-l}\cos(\Omega_0(2n+l))
+(2M-l+1)\cos(\Omega_0 l)
\right].
\]
而
\begin{align*}
    \sum_{n=-M}^{M-l}\cos(\Omega_0(2n+l))&=\frac{1}{\sin\Omega_0}\sum_{n=-M}^{M-l}\cos(\Omega_0(2n+l))\sin\Omega_0\\
    &=\frac{1}{2\sin\Omega_0}\sum_{n=-M}^{M-l}\left[\sin(\Omega_0(2n+l+1))-\sin(\Omega_0(2n+l-1))\right]\\
    &=\frac{\sin\bigl(\Omega_0(2M-l+1)\bigr)}{\sin\Omega_0}
\end{align*}
故
\[
r_M[l]=\frac12\left[
\frac{\sin\bigl(\Omega_0(2M-l+1)\bigr)}{\sin\Omega_0}
+(2M-l+1)\cos(\Omega_0 l)
\right].
\]
同理，当 \(-2M\le l<0\) 时，令 \(|l|=-l\)，可得相同形式的表达式
\[
r_M[l]=\frac12\left[
\frac{\sin\bigl(\Omega_0(2M-|l|+1)\bigr)}{\sin\Omega_0}
+(2M-|l|+1)\cos(\Omega_0 l)
\right].
\]
合并后，对 \(|l|\le 2M\) 有
\[
r_M[l]=\frac12\left[
\frac{\sin\bigl(\Omega_0(2M-|l|+1)\bigr)}{\sin\Omega_0}
+(2M-|l|+1)\cos(\Omega_0 l)
\right].
\]

因此，\(x\) 的自相关函数为
\[
r_x[l]=\lim_{M\to\infty}\frac{1}{2M+1}r_M[l].
\]
代入上式得
\[
r_x[l]=\lim_{M\to\infty}\left[
\frac{\sin\bigl(\Omega_0(2M-|l|+1)\bigr)}{(2M+1)2\sin\Omega_0}
+\frac{2M-|l|+1}{2(2M+1)}\cos(\Omega_0 l)
\right].
\]
第一项有界，除以 \(2M+1\) 后趋于 \(0\)；第二项趋于 \(\dfrac12\cos(\Omega_0 l)\)。于是
\[
r_x[l]=\frac12\cos(\Omega_0 l),\quad l\in\mathbb Z.
\]

其DTFT为
\[
\DTFT(r_x)=\frac{\pi}{2}(\delta_{\Omega_0}+\delta_{-\Omega_0})
\]
因此
\[
S_x(\Omega)=\frac{1}{2\pi}\lim_{M\to\infty}\frac{|X_M(\e^{\im\Omega})|^2}{2M+1}=\frac{1}{2\pi}\DTFT(r_x)=\frac14(\delta_{\Omega_0}+\delta_{-\Omega_0})
\]
\end{example}

基于此，可以看出幅值为1的余弦序列功率为$1/2$，并且可以猜想：对于离散信号也有类似\cref{thm:PSD}的结论。

\begin{theorem}\label{thm:discrete_PSD}
    对于几乎周期的功率序列$x[n]$，设其（作为分布的）傅里叶变换为
    $$X=\sum_{p}a_p\delta_{\Omega_p}$$
    其中$0\le\Omega_k\le \pi$，则其功率谱密度为
    $$S_x=\sum_p |a_p|^2\delta_{\Omega_p}$$
    进一步，设几乎周期的功率序列$y[n]$的傅里叶变换为
    $$Y=\sum_q b_q\delta_{\Omega_q}$$
    则$x[n]$与$y[n]$的互功率谱密度为
    $$S_{xy}=\sum_{k}a_k\overline{b_k}\delta_{\Omega_k}$$
    其中，求和指标$p,q$和$k$可以是有限的，也可以是可数的，因为分布级数由部分和分布的极限定义；$\Omega_k$为两个信号的公共频率，这是因为不同频率处的交叉项在互功率谱密度中的强度为0.
\end{theorem}
\begin{proof}
    由\cref{thm:DTFT_continuous_1}，可以求得：
    $$x[n]=\sum_p a_p\e^{\im\Omega_p n},\quad y[n]=\sum_q b_q\e^{\im\Omega_q n}$$
    代入互相关的定义：
    \begin{align*}
        r_{xy}[l]&=\lim_{M\to\infty}\frac{1}{2M+1}\sum_{n=-M}^M \left(\sum_p a_p\e^{\im\Omega_p (n+l)}\right)\left(\sum_q \overline{b_q}\e^{-\im\Omega_q n}\right)\\
        &=\sum_{p,q}\left(a_p \overline{b_q}\e^{\im\Omega_p l}\cdot\lim_{M\to\infty}\frac{1}{2M+1}\sum_{n=-M}^M \e^{\im(\Omega_p-\Omega_q)n}\right)
    \end{align*}
    其中\begin{align*}
        \sum_{n=-M}^M \e^{\im(\Omega_p-\Omega_q)n}&=\e^{-\im(\Omega_p-\Omega_q)M}\frac{1-\e^{\im(\Omega_p-\Omega_q)(2M+1)}}{1-\e^{\im(\Omega_p-\Omega_q)}}\\
        &=\frac{\e^{-\im(\Omega_p-\Omega_q)(2M+1)/2}-\e^{\im(\Omega_p-\Omega_q)(2M+1)/2}}{\e^{-\im(\Omega_p-\Omega_q)/2}-\e^{\im(\Omega_p-\Omega_q)/2}}\\
        &=\frac{\sin\left((2M+1)(\Omega_p-\Omega_q)/2\right)}{\sin((\Omega_p-\Omega_q)/2)}
    \end{align*}
    因此\begin{align*}
        r_{xy}[l]&=\sum_{p,q}\left(a_p\overline{b_q}\e^{\im\Omega_p l}\cdot\lim_{M\to\infty}\frac{\sin\left((2M+1)(\Omega_p-\Omega_q)/2\right)}{(2M+1)\sin((\Omega_p-\Omega_q)/2)}\right)
    \end{align*}
    其中
    $$\lim_{M\to\infty}\frac{\sin\left((2M+1)(\Omega_p-\Omega_q)/2\right)}{(2M+1)\sin((\Omega_p-\Omega_q)/2)}=\begin{cases}
        1,&\text{if }p=q\\
        0,&\text{if }p\neq q
    \end{cases}$$
    因此
    \begin{align*}
        r_{xy}[l]=\sum_{k}a_k\overline{b_k}\e^{\im\Omega_k l}\tag{$*$}
    \end{align*}
    其中$\Omega_k$为两个信号的公共频率。为了使用\cref{thm:DTFT_continuous_2}得到最终的结论，我们不妨设指标$k$的范围是$\mathbb{N}$，说明部分和序列
    $$\sum_{k=0}^K a_k\overline{b_k}\e^{\im\Omega_k l}$$
    是一致缓增的。利用$(*)$，令$y=x$可得：
    \begin{align*}
        &r_x[l]=\sum_{p} |a_p|^2\e^{\im\Omega_p l}
    \end{align*}
    进而得到：
    \begin{align*}
        P_x&=\lim_{M\to\infty}\frac{1}{2M+1}\sum_{n=-\infty}^\infty |x[n]|^2=r_x[0]=\sum_{p} |a_p|^2
    \end{align*}
    由于$x[n],y[n]$是功率序列，有：
    $$P_x=\sum_{p} |a_p|^2<\infty,\quad P_y=\sum_{q} |b_q|^2<\infty$$
    由柯西-施瓦兹不等式可得：
    \begin{align*}
        \sum_{k=0}^K a_k\overline{b_k}\e^{\im\Omega_k l}&\le\sum_{k=0}^\infty |a_k b_k|\\
        &\le \left(\sum_{k=0}^\infty |a_k|^2\right)^{1/2}\left(\sum_{k=0}^\infty |b_k|^2\right)^{1/2}\\
        &\le \left(\sum_{p} |a_p|^2\right)^{1/2}\left(\sum_{q} |b_q|^2\right)^{1/2}\\
        &=\sqrt{P_x P_y}<\infty
    \end{align*}
    可见$\sum_{k=0}^K a_k\overline{b_k}\e^{\im\Omega_k l}$对$K$有一致上界$\sqrt{P_x P_y}$.

    由\cref{thm:DTFT_continuous_2}及\cref{thm:deiscrete_winner_khinchin}，我们得到：
    $$S_{xy}=\DTFT\left(r_{xy}[l]\right)=\sum_{k}a_k\overline{b_k}\delta_{\Omega_k}$$
\end{proof}